\documentclass{amsart}

\usepackage{keytheorems}
\usepackage{amsthm}
\usepackage{color}

\newkeytheorem{problem}
\newtheorem*{proposition}{Proposition}
\newcommand{\todo}{\textcolor{red}{TO-DO.} }
\newcommand{\temp}{\textcolor{red}{TEMPORARY.} }

\newcommand{\set}[1]{\left\{#1\right\}}
\newcommand{\powset}{{\cal P}}
\newcommand{\img}{_{*}}
\newcommand{\pimg}{^{*}}

\newcommand{\R}{\mathbb{R}}
\newcommand{\Z}{\mathbb{Z}}
\newcommand{\C}{\mathbb{C}}
\newcommand{\N}{\mathbb{N}}
\newcommand{\Q}{\mathbb{Q}}

\numberwithin{equation}{section}

\begin{document}

	\title[Math 251W: Homework \#11] {
		Math 251W: Foundations of Advanced Mathematics \\
		\textmd{Homework \#11: Portfolio Problems \S 6.3 \& 6.4}}
	\author{Carmen Beltran, Brendan Butler, Teddy Wachtler}
	\date{\today}

	\maketitle

	\begin{problem}[manual-num=57]
		\begin{proposition}
			Prove that $7n < 2^n$ for all natural numbers $n$ so that $n \geq 6$.
		\end{proposition}

		\begin{proof}
			We ask whether $7n < 2^n$ for all natural numbers $n$.
			We will show this to be true through induction.

			Take the case where $n = 6$.
			By substitution, we see that $7(6) = 42 < 2^6 = 64$.
			By substitution, we see that $7n < 2^n$ for $n = 6$.

			Let $k$ be an arbitrary natural number. Suppose $k \geq 6$ and $7n < 2^n$ for $n = k$.
			Take the case where $n = k + 1$.
			Since $k \geq 6$, we see that $1 \leq k$.
			By algebra, we see that:
			\begin{align*}
				7 &\leq 7k \\
				7k + 7 &\leq 14k \\
				7(k + 1) &\leq 2 \cdot 7k.
			\end{align*}
			By substitution, we see that $7k < 2^k$.
			By algebra, we see that:
			\begin{align*}
				2 \cdot 7k &< 2 \cdot 2^k \\
				7(k + 1) \leq 2 \cdot 7k &< 2^{k + 1} \\
				7(k + 1) &< 2^{k + 1}.
			\end{align*}
			By substitution, we see that $7n < 2^n$ for $n = k + 1$.

			Therefore, by induction, we see that $7n < 2^n$ for all natural numbers $n$ so that $n \geq 6$.
		\end{proof}
	\end{problem}

	\begin{problem}[manual-num=61]
		\begin{proposition}
			Prove that for all sets $A$ and $B$ and functions $f : A \rightarrow B$, $f \pimg (\bigcup_{i=1}^{n} W_i) =\bigcup_{i=1}^{n} f \pimg (W_i)$ for all natural numbers $n$, where $W_i \subseteq B$ for all natural numbers $i \leq n$.
		\end{proposition}

		\begin{proof}
			Let $A$ and $B$ be arbitrary sets.
			Let $f : A \rightarrow B$ be an arbitrary function.
			Let $n$ be an arbitrary natural number.
			Let $W$ be a family of sets indexed by $i$ so that $W_i \subseteq B$ for all natural numbers $i \leq n$.

			Let $x_0$ be an arbitrary element of $f \pimg (\bigcup_{i=1}^{n} W_i)$.
			By the definition of pre-images, we see that $f(x_0) \in \bigcup_{i=1}^{n} W_i$.
			By the definition of indexed unions, we see that there exists a particular natural number $i_0$ so that $i_0 \leq n$ and $f(x_0) \in W_{i_0}$.
			By the definition of the pre-image, we see that $x_0 \in f \pimg (W_{i_0})$.
			By the definition of indexed unions, we see that $x_0 \in \bigcup_{i=1}^{n} f \pimg (W_i)$.
			By the definition of subsets, since $x_0 \in f \pimg (\bigcup_{i=1}^{n} W_i)$ implies $x_0 \in \bigcup_{i=1}^{n} f \pimg (W_i)$, we see that $f \pimg (\bigcup_{i=1}^{n} W_i) \subseteq \bigcup_{i=1}^{n} f \pimg (W_i)$

			Let $x_1$ be an arbitrary element of $\bigcup_{i=1}^{n} f \pimg (W_i)$.
			By the definition of indexed unions, we see that there exists a particular natural number $i_1$ so that $i_1 \leq n$ and $x_1 \in f \pimg W_{i_1}$.
			By the definition of the pre-image, we see that $f(x_1) \in W_{i_1}$.
			By the definition of indexed unions, we see that $f(x_1) \in \bigcup_{i=1}^{n} W_i$.
			By the definition of the pre-image, we see that $x_1 \in f \pimg (\bigcup_{i=1}^{n} W_i)$.
			By the definition of subsets, since $x_1 \in \bigcup_{i=1}^{n} f \pimg (W_i)$ implies $x_1 \in f \pimg (\bigcup_{i=1}^{n} W_i)$, we see that $\bigcup_{i=1}^{n} f \pimg (W_i) \subseteq f \pimg (\bigcup_{i=1}^{n} W_i)$.

			By the definition of set equality, we see that $f \pimg (\bigcup_{i=1}^{n} W_i) = \bigcup_{i=1}^{n} f \pimg (W_i)$.
			Therefore, for all sets $A$ and $B$ and functions $f : A \rightarrow B$, $f \pimg (\bigcup_{i=1}^{n} W_i) =\bigcup_{i=1}^{n} f \pimg (W_i)$ for all natural numbers $n$, where $W_i \subseteq B$ for all natural numbers $i \leq n$.
		\end{proof}
	\end{problem}

	\begin{problem}[manual-num=62]
		\begin{proposition}
			Prove that $\sum_{i = 1}^{n} F_i = F_{n + 2} - 1$ for all natural numbers $n$.
		\end{proposition}

		\begin{proof}
			We ask whether $\sum_{i = 1}^{n} F_i = F_{n + 2} - 1$ for all natural numbers $n$.
			We will show this to be true through induction.

			Take the case where $n = 1$.
			By the definition of the Fibonacci series and algebra, we see that:
			$\sum_{i = 1}^{1} F_i = F_1 = 1 = 2 - 1 = F_3 - 1$.
			By substitution, we see that $\sum_{i = 1}^{n} F_i = F_{n + 2} - 1$ for $n = 1$.

			Let $k$ be an arbitrary natural number. Suppose $\sum_{i = 1}^{n} F_i = F_{n + 2} - 1$ for $n = k$.
			Take the case where $n = k + 1$.
			By substitution, we see that $\sum_{i = 1}^{k} F_i = F_{k + 2} - 1$.
			By the definition of the Fibonacci series and algebra, we see that:
			\begin{align*}
				\sum_{i = 1}^{k} F_i + F_{k + 1} &= F_{k + 2} - 1 + F_{k + 1} \\
				\sum_{i = 1}^{k + 1} F_i &= F_{k + 3} - 1 \\
				\sum_{i = 1}^{k + 1} F_i &= F_{(k + 1) + 2} - 1.
			\end{align*}
			By substitution, we see that $\sum_{i = 1}^{n} F_i = F_{n + 2} - 1$ for $n = k + 1$.

			Therefore, by induction, we see that $\sum_{i = 1}^{n} F_i = F_{n + 2} - 1$ for all natural numbers $n$.
		\end{proof}

		\begin{proposition}
			Prove that $\sum_{i = 1}^{n} F_i^2 = F_n \cdot F_{n + 1}$ for all natural numbers $n$.
		\end{proposition}

		\begin{proof}
			We ask whether $\sum_{i = 1}^{n} F_i^2 = F_n \cdot F_{n + 1}$ for all natural numbers $n$.
			We will show this to be true through induction.

			Take the case where $n = 1$.
			By the definition of the Fibbonaci series, we see that $\sum_{i = 1}^{1} F_i^2 = F_1^2 = 1^2 = F_1 \cdot F_{1 + 1}$.
			By substitution, we see that $\sum_{i = 1}^{n} F_i^2 = F_n \cdot F_{n + 1}$ for $n = 1$.

			Let $k$ be an arbitrary natural number. Suppose $\sum_{i = 1}^{n} F_i^2 = F_n \cdot F_{n + 1}$ for $n = k$.
			Take the case where $n = k + 1$.
			By substitution, we see that $\sum_{i = 1}^{k} F_i^2 = F_k \cdot F_{k + 1}$.
			\begin{align*}
				\sum_{i = 1}^{k} F_i^2 + F_{k + 1}^2 &= F_k \cdot F_{k + 1} + F_{k + 1}^2 \\
				\sum_{i = 1}^{k + 1} F_i^2 &= F_{k + 1} \cdot (F_k + F_{k + 1})\\
				\sum_{i = 1}^{k + 1} F_i^2 &= F_{k + 1} \cdot (F_{k + 2})\\
				\sum_{i = 1}^{k + 1} F_i^2 &= F_{k + 1} \cdot F_{(k + 1) + 1}.
			\end{align*}
			By substitution, we see that $\sum_{i = 1}^{n} F_i^2 = F_n \cdot F_{n + 1}$ for $n = k + 1$.

			Therefore, by induction, we see that $\sum_{i = 1}^{n} F_i^2 = F_n \cdot F_{n + 1}$ for all natural numbers $n$.
		\end{proof}
	\end{problem}

\end{document}